% ============================================
% LUC-14 (EN)
% English translation — SPEC-935-R355
% ============================================
\section*{\textbf{LUC-14} --- Clinical Self-Assessment (DSM-5-TR)}
\markboth{LUC-14 --- Clinical Self-Assessment (DSM-5-TR)}{}

\noindent \textit{Document extracted from the expert file.}

\vspace{0.5em}\noindent\rule{\textwidth}{0.5pt}\vspace{0.5em}


\noindent {\footnotesize\texttt{PUBLIC ARCHIVE $\cdot$ DOCUMENTARY EVIDENCE $\cdot$ CASE 1,260}}\\[-0.2em]

\noindent\textbf{Document:} unsolicited diagnostic self-assessment.\\
\noindent \textbf{Prepared by:} Dr. Lúcia Mendes --- CRM-MG 28.391 --- Forensic Psychiatrist.\\
\noindent\textbf{Date:} April 13, 2026 (05:47 -- 06:31).\\
\noindent \textbf{Requester:} the subject herself.\\
\noindent \textbf{Purpose:} to verify correspondence between current symptoms and
\noindent DSM-5-TR diagnostic criteria for Brief Psychotic Disorder (298.8),
\noindent Dissociative Identity Disorder (300.14) and Shared Psychotic
\noindent Disorder (Folie à Deux --- 297.3).\\
\noindent \textbf{Document status:} unfinished. The subject interrupted
\noindent completion on page 4.

\noindent \textbf{DIAGNOSTIC CRITERIA --- BRIEF PSYCHOTIC DISORDER (DSM-5-TR 298.8)}

\textbf{A. Presence of one (or more) of the following symptoms.}
\noindent At least one must be (1), (2) or (3).

\noindent \textbf{(1) Delusions.}

\noindent I believe that an entity\slash intelligence\slash pattern
\noindent contaminated a patient in 1917, spread through his diary,
\noindent infected a physician in 1979 and has now contaminated me through reading.

\noindent I recognize, simultaneously, that this belief satisfies
\noindent all the criteria for delusional ideation: it is fixed, idiosyncratic,
\noindent impermeable to counter-evidence, and causes clinically significant distress.

\noindent \textbf{Criterion A(1): present.}

\noindent \textbf{(2) Hallucinations.}

\noindent I report persistent olfactory hallucination (sweet odor,
\noindent sweetish, described as ``wilted flower'' or ``gone-off meat'').

\noindent I report intermittent visual hallucination (perception of
\noindent a yellow glow in the corner of the visual field, with no physical counterpart).

\noindent I report tactile hallucination (sensation of weight on the chest,
\noindent independent of body position, with no identifiable cardiac or respiratory cause).

\noindent I report a localized change in body temperature
\noindent (sensation of heat in the right ocular region, without objective hyperthermia).

\noindent I report a proprioceptive alteration (sensation of a ``hand
\noindent that moves by itself'' during handwriting, with no electromyographic record).

\noindent \textbf{Clinical note:} all the hallucinatory modalities
\noindent reported correspond exactly to those described by the index patient
\noindent (Joaquim Antônio Correia, 1,260) and by the secondary case (Dr. Heitor
\noindent Oliveira). The interindividual coincidence of content, order of
\noindent appearance and temporal progression rules out the hypothesis of
\noindent independent suggestion.

\noindent I consider this clinical note a desperate attempt to
\noindent turn observation into explanation.

\noindent \textbf{Criterion A(2): present.}

\noindent \textbf{(3) Disorganized speech.}

\noindent This document is the proof. A professional with 15 years
\noindent \textit{of experience writing a diagnostic self-assessment} for herself
\noindent at 5:47 in the morning, at home, unsupervised, unsolicited, with
\noindent no legal purpose. If a patient of mine handed me this document,
\noindent I would write in the chart: ``overelaborated speech with persecutory
\noindent content, partial insight, risk of deterioration.''

\noindent \textbf{Criterion A(3): present.}

\noindent \textbf{(4) Grossly disorganized behavior.}

\noindent Not applicable. The subject maintained occupational
\noindent functioning over the last 7 days. She kept appointments. Answered
\noindent phone calls. Replied to e-mails. None of her interlocutors reported
\noindent unusual behavior.

\noindent \textbf{Criterion A(4): absent.}

\noindent \textbf{B. Duration of the episode: at least one day and less than one month.}

\noindent Onset of symptoms: April 5, 2026 (first reading
\noindent of the diary). Current date: April 13, 2026. Duration: 8 days.

\noindent \textbf{Criterion B: present.}

\noindent \textbf{C. Exclusion of affective disorder.}

\noindent No history of major manic or depressive episodes
\noindent in 15 years of clinical practice and 42 years of life. No
\noindent first-degree family history.

\noindent \textbf{Criterion C: present.}

\noindent \textbf{D. Exclusion of substance/medical condition.}

\noindent No use of psychoactive substances. No current medication.
\noindent Laboratory tests requested on April 9 (complete blood count,
\noindent thyroid function, vitamin B12) within normal limits.

\noindent \textbf{Criterion D: present.}

\noindent \textbf{E. Exclusion of another prior psychiatric condition.}

\noindent No prior diagnosis of psychotic, dissociative or
\noindent personality disorder. History of mild generalized anxiety
\noindent (2018--2020), remitted without medication.

\noindent \textbf{Criterion E: present.}

\noindent \textbf{DIFFERENTIAL DIAGNOSIS}

\noindent \textbf{Hypothesis 1: Brief Psychotic Disorder (298.8).}

\noindent All criteria met. It would be the most parsimonious
\noindent diagnosis --- a reactive episode triggered by contact with
\noindent highly traumatic narrative material, with delusional content
\noindent shaped by the material of exposure.

\noindent Counter-indication: the interindividual coincidence of
\noindent symptoms among three people (Joaquim, Oliveira, me) who had no
\noindent contact with one another. Brief Psychotic Episodes are individually
\noindent determined. They do not produce symptomatic replicas in series.

\textbf{Hypothesis 2: Shared Psychotic Disorder}

\noindent Traditionally, it requires direct contact between the
\noindent index patient and the secondary patient(s). Joaquim died in 1979.
\noindent Oliveira died at some point between 1979 and 2026. I never
\noindent met either of them. The transmission occurred exclusively
\noindent through the text.

\noindent There is no precedent in the DSM for folie à deux mediated
\noindent by reading. But there is also no explicit exclusion.

\noindent If I am right, the diary is the first documented
\noindent case of psychotic disorder transmitted by narrative.

\noindent If I am wrong, the diary is a narrative trigger
\noindent powerful enough to produce the same outbreak in three people
\noindent sequentially. A coincidence that 15 years of clinical statistics
\noindent taught me not to accept.

\noindent \textbf{Hypothesis 3: The entity exists.}

\noindent It is not a diagnosis. It is a finding that
\noindent fits neither ICD nor DSM. But clinical observation, when well
\noindent done, does not discard what it cannot classify. It merely records.

\noindent I am recording.

\noindent \textbf{TREATMENT PLAN}

\noindent 1. Return the diary to the archive. Do not open it again.

\noindent 2. Request an independent psychiatric assessment.

\noindent 3. Begin an atypical antipsychotic if symptoms
\noindent persist for more than 14 days.

\noindent 4. Suspend reading of any material related
\noindent to the case.

\noindent 5. Do not write anything by hand anymore.

\noindent 6. Discard this document before anyone finds it.

\noindent 7. Tell no one what is happening.

\noindent \textbf{EVOLUTION (06:31)}

\noindent I interrupted the form 23 minutes ago. I reread what I wrote.
\noindent I identified three errors.

\noindent The first error: the wording is in Joaquim's first voice before
\noindent internment. The sentence structure --- topic + justification + consideration ---
\noindent is the same one he uses in the first pages of the diary.

\noindent I am writing the way he wrote.

\noindent The second error: the treatment plan contradicts
\noindent hypothesis 3. If the entity exists, returning the diary does not interrupt
\noindent the contagion. The contagion is already in me. The diary is only the vector.

\noindent The vector has already fulfilled its function.

\noindent \textit{The third error: I am writing this document against the orientation of my own treatment plan.}

\noindent The plan says: ``Do not write anything by hand anymore.'' The moment I wrote line 7 of the plan, I was already violating line 5.

\noindent The hand wrote before I decided.

\noindent The hand \textit{knew} I was going to continue.

\noindent \textit{The file \texttt{L260-14A.pdf} was found on Dr. Lúcia
\noindent Mendes's computer on May 14, 2026, during the forensic scan of the equipment. The
\noindent document was saved on the desktop, unencrypted, under the name
\texttt{do\_not\_delete\_this.docx}. Forensic analysis confirms that the file was created
\noindent on April 13 and last modified on April 13, at 06:31. There is no
\noindent record of later access. The document was not sent to anyone. The subject
\noindent followed none of the recommendations of her own treatment plan.}




\noindent \textbf{May 2026 --- Belo Horizonte}

\noindent Dr. Lúcia Mendes's last e-mail was sent on May 12, 2026, at 2:47 in the morning.
\noindent The recipient was Dr. Ronaldo Torres, of the Public Archive of Minas Gerais.
\noindent \textquotedblleft{\textit{Ronaldo,}}
\noindent \textit{I need you to keep this account. If something happens to me, I want it on record.}
\noindent \textit{I have read patient 1,260's complete diary. It is 62 years of records --- from 1917 to 1979. The last pages --- according to the record and the gravedigger's account --- would have been written after his death. I know it is impossible, and that is why I do not put a sentence like that in any official report. But the pages are there. The ink dating is inconclusive. The handwriting is the same.}
\noindent \textit{Joaquim Antônio Correia was not insane.}
\noindent \textit{He was possessed.}
\noindent \textit{I know this sounds absurd coming from a forensic psychiatrist. I spent 15 years diagnosing trauma, dissociation, psychosis. I believed everything had a neurological or psychological explanation.}
\noindent \textit{Now I believe there are things psychology does not explain.}
\noindent \textit{I read the diary. Three days later, I started dreaming of the trench. A week later, I smelled it. Two weeks later, my right eye started to lighten.}
\noindent \textit{The same eye. The same yellow.}
\noindent \textit{I had tests done. CT scan. MRI. Blood work. There is no medical explanation.}
\noindent \textit{I went to the Colony. Found a notebook of Joaquim's hidden in a cell --- from the years of mutism. Found an old man who worked there in the 50s. He told me: 'You have read the diary. Now it is your turn.'}
\noindent \textit{I asked my turn for what. He said: 'To host.'}
\noindent \textit{I am afraid, Ronaldo. It is not a rational fear. It is a physical fear. A fear that comes from the stomach. A fear that smells sweet.}
\noindent \textit{Today, I looked in the mirror and saw the yellow eye looking back at me.}
\noindent \textit{It was not my eye.}
\noindent \textit{It was his eye.}
\noindent \textit{Joaquim's eye.}
\noindent \textit{Or of the thing that was inside him.}
\noindent \textit{I will deliver the diary to the archive tomorrow. I will wrap it in a black cloth and leave it at the reception desk. After that, I will disappear for a while.}
\noindent \textit{Do not look for me.}
\noindent \textit{If anyone asks, say the diary has been returned.}
\noindent \textit{And that the cycle continues.}
\noindent Lúcia\textquotedblright{}
\noindent The diary was found at the reception desk of the Public Archive of Minas Gerais on May 13, 2026, wrapped in a black cloth.
\noindent Dr. Lúcia Mendes has not been seen since.
\noindent Her apartment was found empty. The bed was made. The coffee was cold in the cup. The computer was open on a search page: \textquotedblleft{\textit{How to pass the diary on without reading it?}\textquotedblright}
\noindent There was no answer.
\noindent There is no answer.
\noindent The cycle has no reverse.
\noindent The only way out is to go in.
\noindent \textit{Note from the Public Archive of Minas Gerais:}
\noindent \textit{Patient 1,260's diary is currently in our custody.}
\noindent \textit{Waiting for the next reader.}
\noindent \textit{Waiting for you.}
\noindent \textbf{April 2026 --- Belo Horizonte / Quixeramobim}

\noindent After reading Oliveira's letter, Lúcia did what any forensic psychiatrist would do: she followed the trail.

\noindent Dr. Heitor Oliveira disappeared on June 13, 1979. His car was found three weeks later in a parking lot in Quixeramobim, Ceará.

\noindent The gas had run out. The doors were unlocked. Inside, a backpack with clothes, a notebook with 32 written pages (which no one had read), and a bus ticket to Barbacena dated June 17 --- the same day he arrived.

\noindent He had come back.
\noindent Lúcia spent an afternoon at the Public Archive, digging through boxes. She found Oliveira's notebook --- forgotten, dusty, with the label \textquotedblleft{\textit{Belongings of Dr. H. Oliveira --- 1979\textquotedblright}}

\noindent She read the 32 pages in 40 minutes.

\noindent They were clinical notes, but with a difference: the handwriting changed across the pages. It began firm, medical. It ended trembling, irregular, \textit{childish}.
\noindent The last sentence was:
\noindent \textquotedblleft{\textit{I came back to Barbacena. The cycle will not let me go. I return the diary to the archive. I hope the next reader understands better than I did.\textquotedblright{}}}

\noindent \textbf{The witness.}
\noindent The next day, Lúcia found Dona Maria Cândida --- 92 years old, former Colony nurse, now living in Barbacena.
\noindent The interview lasted two hours.
\noindent \textquotedblleft{\textit{Mr. Oliveira? I remember. He arrived in '79, shortly before patient 1,260 died. He was different from the other doctors. He asked questions. He looked people in the eye. The others...\textquotedblright{}} She paused. \textquotedblleft{\textit{The others only applied shock.\textquotedblright{}}}}
\noindent \textquotedblleft{\textit{And after Joaquim's death?\textquotedblright{}}}
\noindent \textquotedblleft{\textit{He changed. He became strange. He spent hours in the ward of the silenced ones, even with no patient there. Once I saw him talking to himself in the empty cell.\textquotedblright{}}}
\noindent \textquotedblleft{\textit{What did he say?\textquotedblright{}}}
\noindent Dona Maria Cândida hesitated.
\noindent \textquotedblleft{\textit{He said: 'I know you are here. I know you hear me. I know you are not a story.'\textquotedblright{}}}
\noindent Lúcia took notes.
\noindent \textquotedblleft{\textit{Anything else?\textquotedblright{}}}
\noindent \textquotedblleft{\textit{On the day he disappeared, he left a note on my desk.\textquotedblright{}}}
\noindent \textquotedblleft{\textit{A note?\textquotedblright{}}}
\noindent \textquotedblleft{\textit{It said: 'Tell the next one. She will come after me.'\textquotedblright{}}}
\noindent Lúcia stopped taking notes.
\noindent \textquotedblleft{\textit{Did you keep the note?\textquotedblright{}}}
\noindent \textquotedblleft{\textit{I kept it.}\textquotedblright} The old woman opened a sewing box. At the bottom, a yellowed paper, folded four times.
\noindent Lúcia reached out her hand. The paper trembled at her fingertips.
\noindent The handwriting was the same as the diary's.

\noindent \textbf{The next reader.}
\noindent Her.
\noindent She was the next one.
\noindent The line had reached her.
\noindent Lúcia spent the night cross-referencing the data she had collected. Oliveira's note, Joaquim's annotations, the autopsy report, the 1979 letter. Every document pointed to the same conclusion --- a conclusion that, as a forensic psychiatrist, she refused to accept.

\noindent At 2:10, she opened three chain-of-custody envelopes.

\noindent In the first, she placed enlarged photographs of Oliveira's handwriting.
\noindent In the second, she scraped a minimal sample of the ink from the note with a sterilized blade.
\noindent In the third, she placed a loose fiber from the black cloth in which the diary had been returned.

\noindent She sealed everything. Signed over the glue. Photographed the seals.

\noindent At 2:47, she sent encrypted copies to a forensic handwriting expert of the Civil Police and to a private laboratory in Contagem. In the subject line, she wrote only: \textit{comparison of historical material --- urgent}.

\noindent Then she put Joaquim's diary in a metal box, closed it with a padlock, and left the key inside a glass of water, on the table.

\noindent It was an absurd gesture. But if she woke up with the dry key in the drawer, she would know she had to stop.

\noindent At 3:14, her phone vibrated.

\noindent The expert had answered too quickly.

\noindent \textquotedblleft{\textit{The handwriting on the note is not Oliveira's. Nor is it a forgery. Call me when you can. Do not read anything else until we talk.\textquotedblright}}

\noindent Lúcia looked at the metal box. The water in the glass was warm. The key was not there.




\vspace{0.5em}
\noindent\textit{\textbf{↪ Routes:} \rota{DOC-01} \rotasep \rota{LUC-02} \rotasep \rota{EPI-01}}
